% organigra1.tex
% Organigrame de la solution - Organisation en trois parties et leurs interactions
% Requires in preamble: \usepackage{tikz}
%   \usetikzlibrary{shapes.geometric, arrows.meta, positioning, fit, backgrounds}
% Colors bleu / bleuu / bleuuu must be defined in preamble.

\begin{tikzpicture}[
    node distance=0.55cm and 2.5cm,
    every node/.style={font=\small},
    % Part headers
    partbox/.style={
        draw, rounded corners=6pt,
        fill=bleu!15, text width=4cm,
        align=center, minimum height=1cm,
        font=\bfseries\small, draw=bleu, line width=1.4pt
    },
    % Sub-items inside each part
    subbox/.style={
        draw, rounded corners=4pt, fill=white,
        text width=3.7cm, align=center,
        minimum height=0.78cm, font=\footnotesize,
        draw=bleuu, line width=0.8pt
    },
    % Sub-items with a lighter tint (secondary components)
    subbox2/.style={
        draw, rounded corners=4pt, fill=bleuuu!6,
        text width=3.7cm, align=center,
        minimum height=0.78cm, font=\footnotesize,
        draw=bleuuu, line width=0.8pt
    },
    % Bidirectional thick relation arrow
    arrow/.style={<->, >=Stealth, thick, draw=bleu},
    % One-way thick arrow
    arrowone/.style={->, >=Stealth, thick, draw=bleu},
    % Lightweight dashed arrow
    dasharrow/.style={->, >=Stealth, dashed, line width=0.8pt, draw=gray!65},
]

% =============================================
% COLUMN 1 : CALCUL
% =============================================
\node[partbox] (calcul) {Partie I --- Calcul};

\node[subbox,  below=0.45cm of calcul] (c1) {Serveurs Dell PowerEdge R720\\(3 noeuds physiques)};
\node[subbox,  below=0.35cm of c1]     (c2) {Cluster Proxmox VE\\Corosync + Quorum};
\node[subbox,  below=0.35cm of c2]     (c3) {Haute Disponibilité (HA)\\basculement automatique des VMs};
\node[subbox2, below=0.35cm of c3]     (c4) {CLBS\\equilibrage de charge proactif\\(CPU + RAM, seuil 50\%)};
\node[subbox,  below=0.35cm of c4]     (c5) {Sona-Web\\HAProxy + Flask (x2) + PostgreSQL};
\node[subbox,  below=0.35cm of c5]     (c6) {Monitoring\\InfluxDB + Prometheus + Grafana};
\node[subbox2, below=0.35cm of c6]     (c7) {Automatisation\\Ansible (3 playbooks) + PXE/TFTP};

% =============================================
% COLUMN 2 : STOCKAGE
% =============================================
\node[partbox, right=2.5cm of calcul] (stockage) {Partie II --- Stockage};

\node[subbox,  below=0.45cm of stockage] (s1) {Cluster Ceph\\3 MON + 2 MGR (actif/standby)};
\node[subbox,  below=0.35cm of s1]       (s2) {3 OSDs (sdb, 300 GB/noeud)\\capacite brute : 900 GB};
\node[subbox2, below=0.35cm of s2]       (s3) {Pool RBD\\replication x3, capacite utile : 300 GB};
\node[subbox,  below=0.35cm of s3]       (s4) {CRUSH Map\\domaine defaillance : host};
\node[subbox2, below=0.35cm of s4]       (s5) {Thin provisioning\\532 GB alloues / 300 GB reel};

% =============================================
% COLUMN 3 : RESEAU
% =============================================
\node[partbox, right=2.5cm of stockage] (reseau) {Partie III --- Reseau};

\node[subbox,  below=0.45cm of reseau] (r1) {Open vSwitch\\bridges OVS sur chaque noeud};
\node[subbox2, below=0.35cm of r1]     (r2) {5 VxLAN d'infrastructure\\Mgmt / Corosync / Ceph / Migration};
\node[subbox,  below=0.35cm of r2]     (r3) {3 reseaux applicatifs\\Dev (60) + Prod (70) + DMZ (80)};
\node[subbox,  below=0.35cm of r3]     (r4) {OPNsense\\routage centralise + NAT + DHCP};
\node[subbox2, below=0.35cm of r4]     (r5) {Pare-feu bi-niveau\\OPNsense (inter-reseaux)\\+ Proxmox Firewall (hyperviseur)};

% =============================================
% GROUPING BACKGROUND BOXES
% =============================================
\begin{scope}[on background layer]
    \node[draw=bleu, rounded corners=9pt, fill=bleu!4,
          line width=1pt, fit=(calcul)(c7), inner sep=7pt] {};
    \node[draw=bleu, rounded corners=9pt, fill=bleu!4,
          line width=1pt, fit=(stockage)(s5), inner sep=7pt] {};
    \node[draw=bleu, rounded corners=9pt, fill=bleu!4,
          line width=1pt, fit=(reseau)(r5), inner sep=7pt] {};
\end{scope}

% =============================================
% INTER-PART RELATIONS (top-level headers)
% =============================================

% Calcul <-> Stockage : volumes RBD pour VMs + live migration CLBS
\draw[arrow] (calcul.east) --
    node[above, font=\tiny\bfseries, text=bleu, yshift=1pt] {volumes RBD (VMs + live migration)}
    (stockage.west);

% Stockage <-> Reseau : trafic de replication Ceph sur VxLAN 30/40
\draw[arrow] (stockage.east) --
    node[above, font=\tiny\bfseries, text=bleu, yshift=1pt] {replication Ceph (VxLAN 30/40)}
    (reseau.west);

% Reseau -> Calcul : segmentation applicative et isolation
\draw[dasharrow, bend left=38] (reseau.north)
    to node[above, font=\tiny, text=gray!80] {segmentation applicative + isolation}
    (calcul.north);

% =============================================
% DETAIL CROSS-RELATIONS (component level)
% =============================================

% CLBS uses InfluxDB metrics AND VxLAN 50 for migration
\draw[dasharrow] (c4.east) to[out=0,in=180, looseness=0.6]
    node[below, font=\scriptsize, text=gray!70, xshift=4pt] {VxLAN 50}
    (r2.west);

% Sona-Web VMs stored on Ceph RBD
\draw[dasharrow] (c5.east) to[out=0,in=180, looseness=0.6]
    node[below, font=\scriptsize, text=gray!70] {disques RBD}
    (s3.west);

% Monitoring Prometheus scrapes Ceph MGR endpoint
\draw[dasharrow] (c6.east) to[out=0,in=180, looseness=0.5]
    node[above, font=\scriptsize, text=gray!70, xshift=6pt] {metriques Ceph}
    (s1.west);

% OPNsense provides DHCP/routing for Ansible provisioning
\draw[dasharrow] (c7.east) to[out=0,in=180, looseness=0.4]
    node[below, font=\scriptsize, text=gray!70, xshift=8pt] {DHCP + routage}
    (r4.west);

\end{tikzpicture}